\documentclass[a4paper, 12pt]{article}

\usepackage{utils}

\renewcommand*{\today}{24 February 2026}

\begin{document}

\hotbox{Algèbre bilinéaire}{CM 3}{\today}

\begin{lemme}{}{}
    Pour toute forme bilinéaire $\varphi$ non identiquement nulle, il existe un vecteur non isotrope par rapport à $\varphi$.
\end{lemme}

\begin{demonstration}
    Supposons que tous les vecteurs soient isotropes.
    Alors pour tout $u, v \in E$, on a
    $$
    \varphi(u, v) = \dfrac{1}{2} \left( \varphi(u + v, u + v) - \varphi(u, u) - \varphi(v, v) \right) = 0 
    $$
    donc $\varphi$ est identiquement nulle, ce qui est une contradiction.
\end{demonstration}

\begin{theoreme}{}{}
    Soit $\varphi$ une forme bilinéaire définie sur un $\K$-espace vectoriel $E$ de dimension fini $n$.
    Alors il existe une base orthogonale (non forcément orthonormale).
\end{theoreme}

\begin{demonstration}
    \begin{enumerate}
        \item 
        Pour $n = 1$, la base formée d'un vecteur non nul est orthogonale.
        \item
        Supposons que le théorème est vrai pour les dimensions inférieures à $n-1$.
        Soit $u$ un vecteur non isotrope.
        $\{u\}^\perp$ est un hyperplan de $E$ de dimension $n-1$.
        \begin{hotwarn}
            À faire avec $E = \{u\}^\perp \oplus Vect(u)$.
        \end{hotwarn}
    \end{enumerate}
\end{demonstration}

\subsection{Produit scalaire}

\begin{definition}{}{}
    $E$ un $\K$-espace vectoriel et $\varphi \in \mathcal{L}_2(E)$ est dit \textbf{un produit scalaire} sur $E$ si
    \begin{itemize}
        \item $\varphi$ est symétrique
        \item $\varphi$ est non dégénérée
        \item $\varphi$ est définie positive, c'est-à-dire que pour tout $u \in E \setminus \{0\}$, $\varphi(u, u) \geq 0$.
    \end{itemize}
\end{definition}

\begin{proposition}{Inégalité de Cauchy-Schwarz}{}
    Soit $E$ un $\K$-espace vectoriel et $\varphi$ un produit scalaire sur $E$.
    Alors on a l'inégalité suivante
    $$
    \forall u, v \in E, \quad \varphi(u, v)^2 \leq \sqrt{\varphi(u, u)} \sqrt{\varphi(v, v)}.
    $$
    avec égalité si $u$ et $v$ sont colinéaires et $\varphi(u, u) = 0 \iff u = 0$.
\end{proposition}

\begin{demonstration}
    \begin{hotwarn}
        À faire avec $P(\lambda) = \varphi(u + \lambda v, u + \lambda v)$ et son discriminant.
    \end{hotwarn}
\end{demonstration}

\begin{proposition}{}{}
    Soit une forme bilinéaire symétrique sur $E$ un $\K$-espace vectoriel est un produit scalaire
    si et seulement si $\varphi$ est définie positive ($\forall u \in E, \varphi(u, u) = 0 \iff u = 0$).
\end{proposition}

\begin{demonstration}
    \begin{description}
        \item[Sens direct.] On suppose que $\varphi$ est un produit scalaire donc $\varphi$ est non dégénérée donc $d_\varphi$ et $g_\varphi$ sont bijectives
        \begin{hotwarn}
            À faire
        \end{hotwarn}
        \item[Sens réciproque.]
    \end{description}
\end{demonstration}

\begin{proposition}{}{}
    Soit $\varphi$ un produit scalaire sur $E$ un $\K$-espace vectoriel.
    Alors tout sous-espace $F$ de $E$ est non isotrope, et de plus, si $dim E = n < +\infty$ alors $E = F \oplus F^\perp$.
\end{proposition}

\begin{demonstration}
    Supposons que $F$ est isotrope, soit $u \in F \cap F^\perp, \varphi(u, u) = 0$, contradiction
    avec le fait que $\varphi$ est un produit scalaire (le seul $u$ qui marche c'est $u = 0$).
    \begin{hotwarn}
        À faire
    \end{hotwarn}
\end{demonstration}

\subsection{Formes quadratiques}

\begin{definition}
    Soit $E$ un $\K$-espace vectoriel.
    Une \textbf{forme quadratique} sur $E$ est une application $q : E \to \K$ telle qu'il existe une forme bilinéaire $\varphi$ sur $E$ vérifiant
    $$
    \forall u \in E, \quad q(u) = \varphi(u, u).
    $$
\end{definition}

\begin{proposition}{}{}
    Soit $q$ une forme quadratique sur $E$ un $\K$-espace vectoriel.
    Alors elle vérifie les propriétés suivantes $\forall (u, v) \in E^2$ et $\forall \lambda \in \K$ :
    \begin{enumerate}
        \item $q(u + v) = q(u) + q(v) + \varphi(u, v) + \varphi(v, u)$
        \item $q(\lambda u) = \lambda^2 q(u)$ (homogénéité de degré 2)
    \end{enumerate}
\end{proposition}

\begin{demonstration}
    \begin{enumerate}
        \item $q(u + v) = \varphi(u + v, u + v) = q(u) + q(v) + \varphi(u, v) + \varphi(v, u)$
        \item $q(\lambda u) = \varphi(\lambda u, \lambda u) = \lambda^2 \varphi(u, u) = \lambda^2 q(u)$
    \end{enumerate}
\end{demonstration}

\begin{definition}
    On note $\mathcal{D}(E)$ l'ensemble des formes quadratiques sur $E$.
\end{definition}

\begin{proposition}{}{}
    Soit $q$ une forme quadratique sur $E$ un $\K$-espace vectoriel.
    Alors
    \begin{enumerate}
        \item $\mathcal{D}(E)$ est un $\K$-espace vectoriel
        \item $q \in \mathcal{D}(E)$ si et seulement si
        \begin{itemize}
            \item $q(\lambda u) = \lambda^2 q(u)$
            \item $\funcdef{\varphi}{E \times E}{\K}{(u, v)}{\dfrac{1}{2} \left( q(u + v) - q(u) - q(v) \right)}$ est une forme bilinéaire symétrique sur $E$
            (On dira que $\varphi$ est la forme polaire de $q$)
        \end{itemize}
    \end{enumerate}
\end{proposition}

\begin{demonstration}
    \begin{hotwarn}
        À faire
    \end{hotwarn}
\end{demonstration}

\begin{corollaire}{}{}
    $\mathcal{S}_2(E)$ est isomorphe à $\mathcal{D}(E)$ et si $dim E = n$ alors $dim \mathcal{D}(E) = \dfrac{n(n + 1)}{2}$.
\end{corollaire}

\begin{demonstration}
    \begin{hotwarn}
        À faire et revoir
    \end{hotwarn}
\end{demonstration}

\subsection{Polynômes homogènes de degré 2}

\begin{proposition}{}{}
    Soit $E$ un $\K$-espace vectoriel de dimension finie $n$ et $\mathcal{B} = (e_1, \dots, e_n)$ une base de $E$.
    Soit $q \in \mathcal{D}(E)$ et $\varphi$ sa forme polaire.

    Alors il existe un unique polynôme $P \in \K_2[X_1, \dots, X_n]$ homogènes (pas de terme constant)
    tel que pour tout $u \in E, u = \sum\limits_{i=1}^n u_i e_i$ on ait
    $$
    q(u) = P(u_1, \dots, u_n)
    $$
    et le polynôme $P$ est donné par
    $$
    P(X_1, \dots, X_n) = \sum\limits_{i=1}^n \varphi(e_i, e_i) X_i^2 + 2 \sum\limits_{1 \leq i < j \leq n} \varphi(e_i, e_j) X_i X_j.
    $$

\end{proposition}

\end{document}